\section{信道模型与一次传输的信息量}
\label{sec:07-Information-Theory/05-Channel-Coding/01}

你对着一个不太可靠的麦克风说话。发``0''，对方以 \(90\%\) 的概率听到``0''，剩下 \(10\%\) 听成``1''；发``1''也对称地翻。线路每秒推过去 1000 个比特。

接收端每秒到手多少信息？

1000 不对，有些位被噪声吃掉了。900 也不对——被吃掉的那一百位不会举手，接收端并不知道该怀疑哪几位。真实答案是 1000 乘上一个小于 1 的折扣因子。要把这个因子算出来，得先把三样东西分开：信道本身、你怎么用它，以及二者合起来诱导出的联合分布。

\subsection{信道是环境给的，输入分布是你手里的自由度}

离散信道就是一张条件概率表

\[
W(y\given x)=P(Y=y\given X=x).
\]

它只回答一个问题：输入既已定为 \(x\)，输出会怎样随机地跳。这是环境的物理属性，发送者改不了；它也没有规定你多久用一次某个输入。

那一部分归你。输入分布 \(p_X\) 说的是你打算以多大频率发 0、多大频率发 1。信道与输入策略凑齐，联合分布才出现：

\[
p_{XY}(x,y)=p_X(x)W(y\given x),\qquad
p_Y(y)=\sum_x p_X(x)W(y\given x).
\]

于是同一张转移表，在不同用法下能承载的信息量不同。极端一点看：你永远只发 0，那么无论线路多干净，接收端也只是一遍遍确认它早就知道的事，信息量为零。信道是死的，输入策略是活的，联合分布是二者相乘的结果——把这三者揉成一团，是后面理解容量时最常见的绊脚石。

\subsection{该量的是输出替输入消掉了多少不确定}

一个顺手的想法是直接看输出熵 \(H(Y)\)：输出变化越丰富，传的信息越多？

反例来得很快。设想信道压根不理会输入，无论你发什么，输出端都以各半的概率吐出 0 和 1。\(H(Y)=1\) 比特（bit），热闹得很，可这串 0/1 与你的输入毫无瓜葛，接收端从中猜不出任何东西。换看输入熵 \(H(X)\) 也不行：信道完美无噪，但你只用一种输入，\(H(X)=0\)，物理上能传的东西你一点没往里装。

两个方向都碰壁，说明该量的既不是发端制造了多少不确定，也不是收端出现了多少变化，而是\emph{看到输出之后，对输入的剩余不确定性缩水了多少}。这正是互信息：

\[
I(X;Y)=H(X)-H(X\given Y)=H(Y)-H(Y\given X).
\]

第二个写法在通信语境里更好用。\(H(Y)\) 是输出端的全部变化，\(H(Y\given X)\) 是输入钉死之后仍在抖动的那一部分——纯噪声。把噪声扣掉，剩下的才是输入能操纵的信号。

还有第三种写法，它把``为什么能传''说得最清楚：

\[
I(X;Y)=\sum_x p_X(x)\,\KL\bigl(W(\cdot\given x)\,\|\,p_Y\bigr).
\]

某个输入 \(x\) 之所以能携带信息，是因为它诱导出的输出分布 \(W(\cdot\given x)\) 与总体输出分布 \(p_Y\) 拉开了距离。若所有输入诱导出同一个输出分布，收端看到的东西对谁都一样，互信息归零。所以关键不在于输出乱不乱，而在于不同输入是否推出可区分的输出。

两个极端可以当刻度用。无噪时 \(Y=X\)，\(H(Y\given X)=0\)，互信息等于 \(H(X)\)——信道原样交付你装进去的东西，装多少交付多少。另一头，若 \(W(y\given x)\) 完全不随 \(x\) 变，则 \(X\) 与 \(Y\) 独立，互信息为零，哪怕 \(H(Y)\) 大得离谱。

\subsection{二元擦除信道：一个能手算的透明例子}

二元擦除信道（BEC）简单到可以口算。它以概率 \(1-\varepsilon\) 原样交付输入比特，以概率 \(\varepsilon\) 交出一个特殊符号 \(?\)，且擦除与输入值无关。

\begin{center}
\begin{tikzpicture}[x=1cm,y=1cm,font=\small,>=stealth]
  \node[font=\footnotesize,black!70] at (0,2.75) {输入 \(X\)};
  \node[font=\footnotesize,black!70] at (4.2,2.75) {输出 \(Y\)};
  \draw[black!55] (0,2) circle (0.3) node {\(0\)};
  \draw[black!55] (0,0) circle (0.3) node {\(1\)};
  \draw[black!55] (4.2,2) circle (0.3) node {\(0\)};
  \draw[black!55] (4.2,1) circle (0.3) node {\(?\)};
  \draw[black!55] (4.2,0) circle (0.3) node {\(1\)};
  \draw[->,blue!65,thick] (0.32,2) -- (3.86,2);
  \draw[->,blue!65,thick] (0.32,0) -- (3.86,0);
  \draw[->,red!60] (0.26,1.85) -- (3.92,1.12);
  \draw[->,red!60] (0.26,0.15) -- (3.92,0.88);
  \node[above,font=\footnotesize,blue!65] at (2.1,2.04) {\(1-\varepsilon\)};
  \node[below,font=\footnotesize,blue!65] at (2.1,-0.04) {\(1-\varepsilon\)};
  \node[font=\footnotesize,red!60] at (2.1,1.68) {\(\varepsilon\)};
  \node[font=\footnotesize,red!60] at (2.1,0.32) {\(\varepsilon\)};
  \node[right,font=\footnotesize,black!60,align=left] at (4.7,1) {收到 \(?\)：知道自己\\什么都不知道};
  \node[right,font=\footnotesize,black!60,align=left] at (4.7,2) {收到 0 或 1：确定};
\end{tikzpicture}

\emph{擦除信道的好处在于它诚实：错误从不伪装成正确。收端要么拿到一个百分之百可信的比特，要么拿到一个明确的``丢了''标记，永远不会被一个看起来正常的翻转位骗过去。}
\end{center}

这份诚实让条件熵一眼可算。收到确定比特时不确定性为零，收到 \(?\) 时后验退回先验、不确定性等于 \(H(X)\)，于是 \(H(X\given Y)=\varepsilon H(X)\)，从而

\[
I(X;Y)=(1-\varepsilon)H(X).
\]

代入数字：公平输入 \(H(X)=1\) 比特、\(\varepsilon=0.2\)，每次使用平均送达 \(0.8\) 比特。但若输入极度偏斜，\(99\%\) 发 0 只有 \(1\%\) 发 1，则 \(H(X)\approx0.08\) 比特，互信息跌到 \(0.064\) 左右。信道纹丝未动，是你装进去的消息太少。

这个对比已经把容量的定义逼出来了：光问``当前用法下传了多少''不够，还得问``最好的用法下能传多少''。

\subsection{无记忆说的是信道，不是消息}

离散无记忆信道（DMC）指的是

\[
P(y^n\given x^n)=\prod_{i=1}^n W(y_i\given x_i).
\]

等号右边按位置分解，含义是第 \(i\) 次的输出在给定第 \(i\) 次输入之后与其余位置无关。这是对噪声说的话。

它并不要求发送出去的符号 \(X_1,\dots,X_n\) 相互独立。恰恰相反，纠错码的码字符号通常高度相关：信息比特经编码器之后，相邻符号之间被刻意塞进了结构性依赖，而这份依赖正是纠错能力的来源——某一位被打坏时，其他位还留着关于它的线索。噪声在各时刻独立作用，不等于你在各时刻不能带着约束发送。

现实里打破这条假设的情形不少：突发错误、随时间漂移的信道状态、反馈、干扰、同步失配。一次雷暴可以让连续几百毫秒的接收信号全部报废，那不是独立噪声，是成片的擦除。转移矩阵是分析的起点，不是所有信道的完整画像。

\subsection{后处理不会凭空生出证据}

若接收端拿到 \(Y\) 之后再做量化、降噪、压缩，得到 \(Z\)，整条链路是 Markov 链 \(X\to Y\to Z\)。数据处理不等式给出

\[
I(X;Z)\le I(X;Y).
\]

\begin{center}
\begin{tikzpicture}[x=1cm,y=1cm,font=\small,>=stealth]
  \draw[black!55] (0,1.4) circle (0.32) node {\(X\)};
  \draw[black!55] (3.2,1.4) circle (0.32) node {\(Y\)};
  \draw[black!55] (6.4,1.4) circle (0.32) node {\(Z\)};
  \draw[->,black!60] (0.34,1.4) -- (2.84,1.4);
  \draw[->,black!60] (3.54,1.4) -- (6.04,1.4);
  \node[font=\footnotesize,black!70] at (1.6,1.72) {物理信道 \(W\)};
  \node[font=\footnotesize,black!70] at (4.8,1.72) {任意后处理};
  \fill[blue!55] (0,0.35) rectangle (4.6,0.65);
  \draw[black!40] (0,0.35) rectangle (4.6,0.65);
  \node[right,font=\footnotesize,black!70] at (4.8,0.5) {\(I(X;Y)\)};
  \fill[blue!55] (0,-0.35) rectangle (2.9,-0.05);
  \draw[black!40] (0,-0.35) rectangle (4.6,-0.05);
  \node[right,font=\footnotesize,black!70] at (4.8,-0.2) {\(I(X;Z)\le I(X;Y)\)};
  \node[font=\footnotesize,black!55,align=center] at (2.3,-0.85) {后处理只会削短这根条，不会加长};
\end{tikzpicture}

\emph{量化、平滑、抽取特征都属于对 \(Y\) 的加工。加工能让信息更易用，却搬不来原本不在 \(Y\) 里的证据。}
\end{center}

于是有一条对工程安排的硬约束：编码增益必须在发送之前布置好。收到一段含噪信号后再聪明的算法也只是在整理已有线索，不能增产。信道编码定理把编码放在发端，原因就在这里——只有事先在输入符号之间造出可区分的结构，收端才有机会把被打乱的消息还原。同一条不等式在数据工作里换了身衣服：特征工程能让下游模型更容易学，但它不会把标签里没有的信息造出来；采集环节丢掉的东西，建模环节补不回来。更完整的讨论见\hyperref[sec:07-Information-Theory/02-Joint-Information/04]{``数据处理不等式与充分表示''}。

\subsection{别急着把互信息叫容量}

\(I(X;Y)\) 依赖于你此刻用的 \(p_X\)。以各半概率发 0 和 1 算出来的那个数，只是这一种用法的成绩单。在所有允许的输入分布上取最大，才得到信道容量

\[
C=\max_{p_X}I(X;Y).
\]

容量只认转移矩阵，不认任何具体策略。互信息像某位车手今天跑出的圈速，容量像这台车在这条赛道上的理论极限——后者不因前者发挥失常而变小。

BEC 那个 \((1-\varepsilon)H(X)\) 里还留着输入分布的尾巴。下一节换成二元对称信道，把``先算任意输入、再优化''这两步完整走一遍，得到本单元最常被引用的一条公式，也顺带回答开头那个麦克风的问题：往 \(10\%\) 翻转率的线路里灌 1000 bit/s，真正过得去的大约是 531 bit/s，其余填的是噪声。
